\documentclass[12pt,a4paper]{article}
\usepackage[utf8]{inputenc}
\usepackage[T1]{fontenc}
\usepackage[french]{babel}
\usepackage{geometry}
\usepackage{fancyhdr}
\usepackage{amsmath}
\usepackage{xcolor}
\usepackage{hyperref}
\usepackage{graphicx}
\usepackage{enumitem}
\usepackage{tcolorbox}

\geometry{margin=2.5cm}
\setlength{\headheight}{14.5pt}

\pagestyle{fancy}
\fancyhf{}
\fancyhead[L]{Résumé Théorique - Administration Unix}
\fancyhead[R]{\thepage}
\fancyfoot[C]{}

\definecolor{infobox}{RGB}{230,240,255}
\newtcolorbox{infobox}[1]{
    colback=infobox,
    colframe=blue!50!black,
    title=#1,
    fonttitle=\bfseries
}

\title{\textbf{Résumé Théorique Détaillé - Administration Unix/Linux}\\
\large Explications Conceptuelles et Mécanismes Internes - Examen Final}
\author{}
\date{}

\begin{document}

\maketitle
\thispagestyle{fancy}

\section*{Introduction}
Ce document présente les explications théoriques détaillées des mécanismes internes et des concepts fondamentaux de l'administration Unix/Linux. Il couvre toutes les parties au programme de l'examen avec des explications approfondies sur le fonctionnement des systèmes.

\textbf{Objectif} : Comprendre le "pourquoi" et le "comment" derrière les commandes et procédures, essentiel pour répondre aux questions théoriques de l'examen.

\newpage
\tableofcontents
\newpage

% ============================================================
% SECTION 1 : HIÉRARCHIE DES SYSTÈMES DE FICHIERS (FHS)
% ============================================================

\section{Hiérarchie des Systèmes de Fichiers (FHS) - Théorie}

\subsection{Principe du FHS}

Le \textbf{Filesystem Hierarchy Standard (FHS)} est un standard qui définit l'organisation des répertoires dans les systèmes Unix/Linux. Il garantit :

\begin{itemize}
    \item \textbf{Cohérence} : Même organisation entre différentes distributions
    \item \textbf{Portabilité} : Facilite le passage d'une distribution à l'autre
    \item \textbf{Maintenabilité} : Localisation prévisible des fichiers
    \item \textbf{Séparation des préoccupations} : Distinction entre fichiers système et utilisateur
\end{itemize}

\subsection{Architecture des répertoires}

\subsubsection{Distinction /bin vs /usr/bin}

\begin{infobox}{Concept Important}
La distinction entre /bin et /usr/bin reflète une séparation historique et fonctionnelle :
\end{infobox}

\begin{itemize}
    \item \textbf{/bin et /sbin} : Dans la partition racine (/)
    \begin{itemize}
        \item Nécessaires au démarrage minimal (mode single-user)
        \item Disponibles avant montage de /usr
        \item Contiennent les commandes essentielles
    \end{itemize}
    \item \textbf{/usr/bin et /usr/sbin} : Peuvent être sur partition séparée
    \begin{itemize}
        \item Montés après la racine
        \item Contiennent les programmes additionnels
        \item Peuvent être en lecture seule
    \end{itemize}
\end{itemize}

\subsubsection{Organisation logique}

Le FHS organise les répertoires selon leur fonction :

\begin{enumerate}
    \item \textbf{Répertoires statiques} : Contenu ne change pas pendant fonctionnement
    \begin{itemize}
        \item /bin, /sbin, /lib : Exécutables et bibliothèques système
        \item /usr : Programmes et données utilisateur
        \item /opt : Logiciels optionnels
    \end{itemize}
    \item \textbf{Répertoires variables} : Contenu change pendant fonctionnement
    \begin{itemize}
        \item /var : Logs, caches, files d'attente
        \item /tmp : Fichiers temporaires
        \item /run : Données runtime (PID, sockets)
    \end{itemize}
    \item \textbf{Répertoires utilisateur} : Données utilisateur
    \begin{itemize}
        \item /home : Répertoires personnels
        \item /root : Répertoire de l'administrateur
    \end{itemize}
\end{enumerate}

\subsection{Types de fichiers Linux - Théorie}

\subsubsection{Inodes : Concept fondamental}

\begin{infobox}{Définition}
Un \textbf{inode} (index node) est une structure de données contenant toutes les métadonnées d'un fichier, sauf son nom.
\end{infobox}

\paragraph{Contenu d'un inode}
\begin{itemize}
    \item \textbf{Type} : Fichier régulier, répertoire, lien symbolique, périphérique, etc.
    \item \textbf{Permissions} : Droits d'accès (rwx) pour propriétaire, groupe, autres
    \item \textbf{Propriétaire} : UID et GID
    \item \textbf{Taille} : Taille en octets
    \item \textbf{Dates} : atime (accès), mtime (modification), ctime (changement métadonnées)
    \item \textbf{Pointeurs vers blocs} : Adresses des blocs de données sur disque
    \item \textbf{Compteur de références} : Nombre de liens physiques
\end{itemize}

\paragraph{Important : Nom du fichier}
\begin{itemize}
    \item Le \textbf{nom du fichier n'est PAS dans l'inode}
    \item Il est stocké dans le \textbf{répertoire parent}
    \item Un répertoire est un fichier spécial contenant des associations nom-inode
    \item Plusieurs noms peuvent pointer vers le même inode (liens physiques)
\end{itemize}

\subsubsection{Liens : Mécanisme de référence}

\paragraph{Liens physiques (Hard Links)}
\begin{itemize}
    \item \textbf{Principe} : Plusieurs entrées de répertoire pointent vers le même inode
    \item \textbf{Compteur} : L'inode maintient un compteur de références
    \item \textbf{Suppression} : Le fichier n'est supprimé que quand le compteur atteint 0
    \item \textbf{Limitation} : Doit rester dans le même système de fichiers
    \item \textbf{Avantage} : Pas de risque de lien cassé
\end{itemize}

\paragraph{Liens symboliques (Soft Links)}
\begin{itemize}
    \item \textbf{Principe} : Fichier spécial contenant le chemin vers la cible
    \item \textbf{Inode séparé} : Le lien a son propre inode
    \item \textbf{Contenu} : Le chemin (absolu ou relatif) vers le fichier cible
    \item \textbf{Avantage} : Peut traverser les systèmes de fichiers
    \item \textbf{Inconvénient} : Devient "cassé" si la cible est supprimée
\end{itemize}

\subsubsection{Types de fichiers spéciaux}

\begin{itemize}
    \item \textbf{Périphériques caractère (c)} : Accès séquentiel (terminaux, ports série)
    \item \textbf{Périphériques bloc (b)} : Accès par blocs (disques)
    \item \textbf{Pipes nommés (FIFO) (p)} : Communication inter-processus
    \item \textbf{Sockets (s)} : Communication réseau locale
\end{itemize}

% ============================================================
% SECTION 2 : DÉMARRAGE DU SYSTÈME - THÉORIE
% ============================================================

\section{Démarrage du Système - Théorie Détaillée}

\subsection{Séquence de démarrage : Mécanismes internes}

\subsubsection{Phase 1 : Matériel et Firmware}

\begin{enumerate}
    \item \textbf{Power-on} : Alimentation du système
    \item \textbf{POST (Power-On Self-Test)} : Vérification matérielle de base
    \item \textbf{Initialisation matérielle} : Détection CPU, RAM, périphériques
    \item \textbf{Chargement du firmware} : BIOS ou UEFI prend le contrôle
\end{enumerate}

\subsubsection{Phase 2 : Bootloader (GRUB)}

\begin{infobox}{Rôle du Bootloader}
Le bootloader est le premier logiciel exécuté après le firmware. Il a pour mission de charger le noyau Linux en mémoire.
\end{infobox}

\paragraph{Processus de chargement GRUB}
\begin{enumerate}
    \item \textbf{MBR/GPT} : Le firmware lit le secteur 0 (MBR) ou la partition EFI
    \item \textbf{GRUB Stage 1} : Code minimal dans MBR (512 octets)
    \item \textbf{GRUB Stage 2} : Code complet chargé depuis /boot/grub2/
    \item \textbf{Menu} : Affichage du menu de sélection (si configuré)
    \item \textbf{Chargement noyau} : Lecture du noyau et initramfs depuis disque
    \item \textbf{Passage au noyau} : Transfert du contrôle avec paramètres
\end{enumerate}

\paragraph{Initramfs : Pourquoi ?}
\begin{itemize}
    \item \textbf{Problème} : Le noyau a besoin de modules pour accéder au disque racine
    \item \textbf{Solution} : Initramfs contient ces modules en RAM
    \item \textbf{Scénarios nécessaires} :
    \begin{itemize}
        \item Disque racine sur LVM (Logical Volume Manager)
        \item Disque racine sur RAID
        \item Système de fichiers crypté
        \item Système de fichiers réseau (NFS)
    \end{itemize}
    \item \textbf{Processus} : Initramfs monte la racine, puis passe le contrôle
\end{itemize}

\subsubsection{Phase 3 : Noyau Linux}

\paragraph{Initialisation du noyau}
\begin{enumerate}
    \item \textbf{Initialisation interne} : Structures de données du noyau
    \item \textbf{Détection matérielle} : CPU, mémoire, périphériques
    \item \textbf{Chargement modules} : Pilotes depuis initramfs ou /lib/modules
    \item \textbf{Montage racine} : Accès au système de fichiers racine
    \item \textbf{Libération initramfs} : Mémoire libérée pour le système
\end{enumerate}

\paragraph{Processus PID 1}
\begin{itemize}
    \item Le noyau lance le premier processus utilisateur (PID = 1)
    \item Historiquement : /sbin/init (SysVinit)
    \item Moderne : systemd (remplace init)
    \item Ce processus est le parent de tous les autres processus
\end{itemize}

\subsection{GRUB2 : Architecture et Configuration}

\subsubsection{Architecture GRUB2}

\begin{itemize}
    \item \textbf{Modulaire} : GRUB2 est modulaire (contrairement à GRUB Legacy)
    \item \textbf{Scripts} : /etc/grub.d/ contient des scripts de génération
    \item \textbf{Configuration} : /etc/default/grub contient les paramètres
    \item \textbf{Génération} : grub2-mkconfig génère grub.cfg depuis ces sources
\end{itemize}

\subsubsection{Pourquoi ne pas éditer grub.cfg directement ?}

\begin{infobox}{Important}
Le fichier /boot/grub2/grub.cfg est généré automatiquement et ne doit JAMAIS être édité directement.
\end{infobox}

\begin{itemize}
    \item \textbf{Raison 1} : Régénération écrase les modifications
    \item \textbf{Raison 2} : Syntaxe complexe, risque d'erreur
    \item \textbf{Raison 3} : Scripts grub.d/ peuvent être mis à jour
    \item \textbf{Solution} : Modifier /etc/default/grub puis régénérer
\end{itemize}

\subsection{Systemd : Révolution du démarrage}

\subsubsection{Problèmes de SysVinit}

\begin{itemize}
    \item \textbf{Séquentiel} : Démarrage séquentiel (lent)
    \item \textbf{Scripts shell} : Scripts complexes, difficiles à déboguer
    \item \textbf{Pas de dépendances} : Ordre basé sur numérotation
    \item \textbf{Logs dispersés} : Logs dans /var/log, format varié
\end{itemize}

\subsubsection{Avantages de systemd}

\begin{itemize}
    \item \textbf{Parallélisation} : Démarrage parallèle des services indépendants
    \item \textbf{Dépendances} : Gestion automatique des dépendances
    \item \textbf{Unités} : Configuration déclarative (fichiers .service)
    \item \textbf{Journal intégré} : journalctl pour tous les logs
    \item \textbf{Activation à la demande} : Services démarrés quand nécessaires
\end{itemize}

\subsubsection{Concepts systemd}

\paragraph{Unités}
\begin{itemize}
    \item \textbf{Service} : Service système (httpd.service)
    \item \textbf{Target} : Groupe d'unités (graphical.target)
    \item \textbf{Mount} : Point de montage (home.mount)
    \item \textbf{Socket} : Socket d'écoute (sshd.socket)
    \item \textbf{Timer} : Tâche planifiée (backup.timer)
\end{itemize}

\paragraph{Dépendances}
\begin{itemize}
    \item \textbf{Requires} : Dépendance forte (si échec, arrêt)
    \item \textbf{Wants} : Dépendance faible (si échec, continue)
    \item \textbf{After} : Ordre de démarrage
    \item \textbf{Before} : Ordre inverse
\end{itemize}

% ============================================================
% SECTION 3 : GESTION DES PAQUETAGES - THÉORIE
% ============================================================

\section{Gestion des Paquetages - Théorie}

\subsection{Concept de paquetage}

\subsubsection{Structure d'un paquet}

Un paquet est une archive contenant :
\begin{itemize}
    \item \textbf{Fichiers} : Binaires, bibliothèques, configuration, documentation
    \item \textbf{Métadonnées} : Nom, version, description, dépendances
    \item \textbf{Scripts} : Pre-install, post-install, pre-remove, post-remove
    \item \textbf{Signature} : Vérification d'intégrité et authenticité
\end{itemize}

\subsubsection{Gestion des dépendances}

\begin{infobox}{Problème des dépendances}
Un logiciel A peut nécessiter des bibliothèques B et C. Comment gérer cela automatiquement ?
\end{infobox}

\paragraph{Résolution automatique}
\begin{itemize}
    \item \textbf{Dépendances directes} : Paquets explicitement requis
    \item \textbf{Dépendances transitives} : Dépendances des dépendances
    \item \textbf{Conflits} : Paquets incompatibles détectés
    \item \textbf{Base de données} : Métadonnées stockées dans /var/lib/rpm/
\end{itemize}

\paragraph{Avantages}
\begin{itemize}
    \item Installation simplifiée : Une commande installe tout
    \item Cohérence : Versions compatibles garanties
    \item Sécurité : Mises à jour de sécurité propagées
\end{itemize}

\subsection{RPM : Format et mécanismes}

\subsubsection{Structure RPM}

\begin{itemize}
    \item \textbf{Format} : Archive cpio compressée
    \item \textbf{En-tête} : Métadonnées (nom, version, dépendances)
    \item \textbf{Payload} : Fichiers à installer
    \item \textbf{Signature} : Vérification cryptographique
\end{itemize}

\subsubsection{Base de données RPM}

\begin{itemize}
    \item \textbf{Localisation} : /var/lib/rpm/
    \item \textbf{Contenu} : Métadonnées de tous les paquets installés
    \item \textbf{Requêtes} : rpm -q interroge cette base
    \item \textbf{Vérification} : rpm -V compare fichiers avec base
\end{itemize}

\subsection{YUM/DNF : Gestionnaires haut niveau}

\subsubsection{Architecture YUM}

\begin{itemize}
    \item \textbf{Dépôts} : Sources de paquets (HTTP, FTP, local)
    \item \textbf{Métadonnées} : Cache des informations des dépôts
    \item \textbf{Résolution} : Algorithme de résolution de dépendances
    \item \textbf{Transactions} : Installation/Suppression atomique
\end{itemize}

\subsubsection{Processus d'installation}

\begin{enumerate}
    \item \textbf{Requête} : \texttt{yum install package}
    \item \textbf{Recherche} : YUM cherche dans les dépôts configurés
    \item \textbf{Résolution} : Calcul des dépendances nécessaires
    \item \textbf{Téléchargement} : Téléchargement des paquets requis
    \item \textbf{Vérification} : Vérification des signatures
    \item \textbf{Installation} : Exécution des scripts et copie des fichiers
    \item \textbf{Mise à jour base} : Mise à jour de /var/lib/rpm/
\end{enumerate}

% ============================================================
% SECTION 4 : GESTION DES UTILISATEURS - THÉORIE
% ============================================================

\section{Gestion des Utilisateurs et Groupes - Théorie}

\subsection{Modèle de sécurité Unix}

\subsubsection{Principe fondamental}

\begin{infobox}{Principe de moindre privilège}
Chaque processus s'exécute avec les privilèges de l'utilisateur qui l'a lancé, limitant les dommages en cas de compromission.
\end{infobox}

\subsubsection{Identifiants}

\begin{itemize}
    \item \textbf{UID (User ID)} : Identifiant unique utilisateur
    \begin{itemize}
        \item 0 : root (super-utilisateur)
        \item 1-999 : Comptes système
        \item 1000+ : Utilisateurs normaux
    \end{itemize}
    \item \textbf{GID (Group ID)} : Identifiant de groupe
    \item \textbf{Groupes} : Un utilisateur peut appartenir à plusieurs groupes
    \begin{itemize}
        \item Groupe primaire : Défini dans /etc/passwd
        \item Groupes secondaires : Définis dans /etc/group
    \end{itemize}
\end{itemize}

\subsection{Fichiers de configuration : Structure et sécurité}

\subsubsection{/etc/passwd : Historique et évolution}

\begin{itemize}
    \item \textbf{Historique} : Contenait autrefois les mots de passe cryptés
    \item \textbf{Problème} : Fichier lisible par tous (permissions 644)
    \item \textbf{Solution} : Mots de passe déplacés dans /etc/shadow
    \item \textbf{Champ passwd} : Contient maintenant "x" (indique présence de shadow)
\end{itemize}

\subsubsection{/etc/shadow : Sécurité renforcée}

\begin{itemize}
    \item \textbf{Permissions} : 000 ou 400 (root uniquement)
    \item \textbf{Hash} : Mots de passe cryptés (MD5, SHA-256, SHA-512, bcrypt)
    \item \textbf{Politiques} : Durées de validité, expiration
    \item \textbf{Sécurité} : Protection contre attaques par dictionnaire
\end{itemize}

\paragraph{Champs de /etc/shadow}
\begin{itemize}
    \item \textbf{last} : Jours depuis 1/1/1970 (epoch)
    \item \textbf{may} : Jours minimum avant changement (empêche changement trop fréquent)
    \item \textbf{must} : Jours maximum avant expiration (force changement)
    \item \textbf{warn} : Jours d'avertissement avant expiration
    \item \textbf{inactif} : Jours d'inactivité avant désactivation
    \item \textbf{disable} : Date d'expiration absolue du compte
\end{itemize}

\subsection{Permissions : Modèle discrétionnaire}

\subsubsection{Modèle DAC (Discretionary Access Control)}

\begin{itemize}
    \item \textbf{Discrétionnaire} : Le propriétaire décide des permissions
    \item \textbf{Trois catégories} : Propriétaire (u), Groupe (g), Autres (o)
    \item \textbf{Trois droits} : Lecture (r), Écriture (w), Exécution (x)
    \item \textbf{Calcul octal} : r=4, w=2, x=1 (somme = permission)
\end{itemize}

\subsubsection{Droits spéciaux}

\paragraph{Setuid (Set User ID)}
\begin{itemize}
    \item \textbf{Principe} : Programme exécuté avec privilèges du propriétaire
    \item \textbf{Exemple} : /usr/bin/passwd (modifie /etc/shadow, propriété root)
    \item \textbf{Risque} : Compromission = élévation de privilèges
    \item \textbf{Utilisation} : Limiter aux programmes nécessitant vraiment root
\end{itemize}

\paragraph{Setgid (Set Group ID)}
\begin{itemize}
    \item \textbf{Sur fichiers} : Exécution avec privilèges du groupe
    \item \textbf{Sur répertoires} : Nouveaux fichiers héritent du groupe du répertoire
    \item \textbf{Utilité} : Collaboration en équipe (fichiers créés appartiennent au groupe)
\end{itemize}

\paragraph{Sticky Bit}
\begin{itemize}
    \item \textbf{Sur répertoires} : Empêche suppression par non-propriétaires
    \item \textbf{Exemple} : /tmp (tous peuvent créer, seul propriétaire peut supprimer)
    \item \textbf{Sécurité} : Protection contre suppression accidentelle
\end{itemize}

\subsection{Sudo : Élévation de privilèges}

\subsubsection{Principe}

\begin{itemize}
    \item \textbf{Problème} : Besoin de privilèges root pour certaines tâches
    \item \textbf{Solution 1} : Connexion en tant que root (risqué)
    \item \textbf{Solution 2} : Sudo (substitute user do)
    \begin{itemize}
        \item Exécution temporaire avec privilèges élevés
        \item Traçabilité : Logs de toutes les commandes sudo
        \item Granularité : Permissions spécifiques par utilisateur/commande
    \end{itemize}
\end{itemize}

\subsubsection{Configuration /etc/sudoers}

\paragraph{Format}
\texttt{utilisateur hôte=(utilisateur\_cible) commandes [NOPASSWD:]}

\paragraph{Exemples}
\begin{itemize}
    \item \texttt{jean ALL=(ALL:ALL) ALL} : Privilèges complets partout
    \item \texttt{\%wheel ALL=(ALL) ALL} : Groupe wheel avec privilèges
    \item \texttt{marie server1=(root) NOPASSWD:/usr/sbin/reboot} : Commande spécifique sans mot de passe
\end{itemize}

\paragraph{Sécurité}
\begin{itemize}
    \item \textbf{visudo} : Éditeur sécurisé (vérifie syntaxe avant sauvegarde)
    \item \textbf{Verrouillage} : Évite édition simultanée
    \item \textbf{Validation} : Rejette syntaxe incorrecte
\end{itemize}

% ============================================================
% SECTION 5 : GESTION DES DISQUES - THÉORIE
% ============================================================

\section{Gestion des Disques - Théorie}

\subsection{Partitionnement : Concepts}

\subsubsection{MBR vs GPT : Différences fondamentales}

\begin{infobox}{Limitations historiques}
Le MBR date de 1983 et reflète les limitations de l'époque (disques < 2 To, 4 partitions).
\end{infobox}

\paragraph{MBR (Master Boot Record)}
\begin{itemize}
    \item \textbf{Localisation} : Secteur 0 du disque (512 octets)
    \item \textbf{Structure} : Table de partitions (4 entrées de 16 octets)
    \item \textbf{Limitations} :
    \begin{itemize}
        \item 4 partitions primaires maximum
        \item Partition étendue pour contourner (logiques)
        \item Adressage 32 bits : 2.2 To maximum par partition
    \end{itemize}
    \item \textbf{Compatibilité} : BIOS uniquement
\end{itemize}

\paragraph{GPT (GUID Partition Table)}
\begin{itemize}
    \item \textbf{Localisation} : Plusieurs secteurs (redondance)
    \item \textbf{Structure} : Table extensible (128 partitions standard, extensible)
    \item \textbf{Avantages} :
    \begin{itemize}
        \item Adressage 64 bits : Jusqu'à 9.4 Zo (zettabytes)
        \item Redondance : Copie de la table
        \item CRC32 : Vérification d'intégrité
    \end{itemize}
    \item \textbf{Compatibilité} : UEFI natif, BIOS avec support
\end{itemize}

\subsection{Systèmes de fichiers : Organisation}

\subsubsection{Structure d'un système de fichiers}

\begin{enumerate}
    \item \textbf{Superblock} : Métadonnées du système de fichiers
    \begin{itemize}
        \item Taille, nombre d'inodes, nombre de blocs
        \item Type de système de fichiers
        \item État (propre ou nécessite fsck)
    \end{itemize}
    \item \textbf{Inode table} : Table des inodes
    \item \textbf{Data blocks} : Blocs de données des fichiers
    \item \textbf{Block bitmap} : Carte des blocs libres/occupés
    \item \textbf{Inode bitmap} : Carte des inodes libres/occupés
\end{enumerate}

\subsubsection{Journalisation : Principe}

\begin{infobox}{Problème sans journalisation}
Sans journalisation, un crash pendant écriture peut corrompre le système de fichiers, nécessitant un fsck long.
\end{infobox}

\paragraph{Principe de la journalisation}
\begin{enumerate}
    \item \textbf{Journal} : Zone spéciale du disque
    \item \textbf{Écriture} : Modifications d'abord écrites dans le journal
    \item \textbf{Commit} : Puis appliquées aux données réelles
    \item \textbf{Récupération} : En cas de crash, rejouer le journal
\end{enumerate}

\paragraph{Types de journalisation}
\begin{itemize}
    \item \textbf{Journalisation métadonnées} : Seulement métadonnées (ext3 par défaut)
    \item \textbf{Journalisation complète} : Métadonnées + données (plus sûr, plus lent)
    \item \textbf{Mode ordonné} : Données écrites avant métadonnées
\end{itemize}

\subsubsection{Extents : Optimisation ext4}

\begin{itemize}
    \item \textbf{Problème} : Fichiers fragmentés = nombreux pointeurs
    \item \textbf{Solution} : Extents (plages contiguës de blocs)
    \item \textbf{Avantage} : Moins de pointeurs, meilleures performances
    \item \textbf{Exemple} : Fichier de 1 Go = 1 extent au lieu de 256K pointeurs
\end{itemize}

\subsection{Montage : Mécanismes}

\subsubsection{Concept de montage}

\begin{infobox}{Abstraction Unix}
Unix présente tous les systèmes de fichiers comme une arborescence unique, même s'ils sont sur différents disques.
\end{infobox}

\paragraph{Processus de montage}
\begin{enumerate}
    \item \textbf{Vérification} : Le système de fichiers existe-t-il ?
    \item \textbf{Validation} : Vérification de l'intégrité (option fsck)
    \item \textbf{Point de montage} : Répertoire où attacher
    \item \textbf{Intégration} : Ajout à l'arborescence unique
    \item \textbf{Cache} : Mise en cache des métadonnées
\end{enumerate}

\subsubsection{/etc/fstab : Montage automatique}

\begin{itemize}
    \item \textbf{Objectif} : Montage automatique au démarrage
    \item \textbf{Exécution} : \texttt{mount -a} lit /etc/fstab
    \item \textbf{UUID} : Préféré à /dev/sdX (plus stable)
    \item \textbf{Options} : Contrôlent le comportement du montage
\end{itemize}

% ============================================================
% SECTION 6 : GESTION DU SWAP - THÉORIE
% ============================================================

\section{Gestion de l'Espace Swap - Théorie}

\subsection{Mémoire virtuelle : Concept}

\subsubsection{Problème}

\begin{itemize}
    \item \textbf{RAM limitée} : Quantité de mémoire physique limitée
    \item \textbf{Besoin} : Exécuter plus de processus que la RAM ne le permet
    \item \textbf{Solution} : Utiliser l'espace disque comme extension de RAM
\end{itemize}

\subsubsection{Mécanisme de pagination}

\begin{enumerate}
    \item \textbf{Pages} : Mémoire divisée en pages (généralement 4 Ko)
    \item \textbf{Pages actives} : En RAM
    \item \textbf{Pages inactives} : Peuvent être déplacées vers swap
    \item \textbf{Page fault} : Accès à une page en swap déclenche chargement
    \item \textbf{Algorithme LRU} : Least Recently Used (pages peu utilisées swapées)
\end{enumerate}

\subsection{Swap : Types et configuration}

\subsubsection{Types de swap}

\begin{itemize}
    \item \textbf{Partition swap} : Partition dédiée
    \begin{itemize}
        \item Performance : Meilleure (accès direct)
        \item Taille : Fixe (nécessite redimensionnement partition)
    \end{itemize}
    \item \textbf{Fichier swap} : Fichier sur système de fichiers
    \begin{itemize}
        \item Flexibilité : Taille modifiable facilement
        \item Performance : Légèrement inférieure (via système de fichiers)
    \end{itemize}
\end{itemize}

\subsubsection{Performance et optimisation}

\begin{itemize}
    \item \textbf{Swappiness} : Paramètre contrôlant l'agressivité du swap
    \begin{itemize}
        \item 0 : Swap seulement si nécessaire
        \item 100 : Swap agressif
        \item Défaut : 60 (généralement)
    \end{itemize}
    \item \textbf{Commande} : \texttt{cat /proc/sys/vm/swappiness}
    \item \textbf{Modification} : \texttt{sudo sysctl vm.swappiness=10}
    \item \textbf{Thrashing} : Trop de swap = système très lent
\end{itemize}

% ============================================================
% CONCLUSION
% ============================================================

\section*{Conclusion}

Ce résumé théorique a présenté les mécanismes internes et concepts fondamentaux de l'administration Unix/Linux :

\begin{itemize}
    \item \textbf{Architecture} : Organisation en couches, séparation noyau/utilisateur
    \item \textbf{FHS} : Organisation logique des répertoires
    \item \textbf{Inodes} : Mécanisme fondamental de gestion des fichiers
    \item \textbf{Démarrage} : Séquence complète du boot, rôle de GRUB et systemd
    \item \textbf{Paquetages} : Gestion des dépendances et résolution automatique
    \item \textbf{Sécurité} : Modèle de permissions, sudo, fichiers shadow
    \item \textbf{Disques} : Partitionnement, systèmes de fichiers, journalisation
    \item \textbf{Swap} : Mémoire virtuelle et pagination
\end{itemize}

Comprendre ces concepts théoriques est essentiel pour :
\begin{itemize}
    \item Répondre aux questions de cours de l'examen
    \item Expliquer les procédures d'administration
    \item Diagnostiquer les problèmes système
    \item Prendre des décisions éclairées en administration
\end{itemize}

\end{document}

